% Задача Коши на кривой: характеристики пересекают gamma трансверсально.
\begin{figure}[htbp]
\centering
\begin{tikzpicture}[>={Latex[length=4pt]}]
  \draw[->] (-0.3,0) -- (6.3,0) node[anchor=north east] {$x$};
  \draw[->] (0,-0.3) -- (0,3.4) node[anchor=south east] {$t$};

  % характеристики x - at = const
  \foreach \c in {-0.6, 0.4, 1.4, 2.4, 3.4, 4.4} {
    \draw[thick, blue!55!black] (\c,0) -- ({\c+1.9},3.05);
  }

  % кривая gamma
  \draw[very thick, red!70!black]
    plot[variable=\s, domain=0.2:5.6, samples=40, smooth]
    (\s, {0.35 + 1.5/(1+exp(-1.3*(\s-2.6)))});

  \node[red!70!black, anchor=west, font=\footnotesize] at (5.0,0.55)
    {$\gamma = \{(\xi(\tau),\, \eta(\tau))\}$};
  \node[red!70!black, anchor=south east, font=\footnotesize] at (2.0,0.62)
    {$u|_\gamma = \varphi(\tau)$};

  \draw[->, red!70!black, very thick] (2.6,1.1) -- (3.35,1.62);
  \node[red!70!black, anchor=north west, font=\footnotesize] at (3.2,1.6)
    {$(\xi',\eta')$};
  \draw[->, blue!55!black, very thick] (2.6,1.1) -- (2.95,1.66);
  \node[blue!55!black, anchor=south east, font=\footnotesize] at (2.72,1.72)
    {$\bar a$};
  \fill (2.6,1.1) circle (1.5pt);
\end{tikzpicture}
\caption{Задача Коши на кривой $\gamma$. Условие трансверсальности требует, чтобы
направляющий вектор $\bar a$ нигде не был касательным к $\gamma$: тогда каждая
характеристика пересекает $\gamma$ ровно один раз и переносит с неё значение
$\varphi$.}
\label{fig:cauchy-gamma}
\end{figure}
